\section{Shared definitions}

\paragraph{Conventions.} These fix every symbol used below; the paper's \S3.1 is identical.
\begin{itemize}
\item $\N:=\{0,1,2,\dots\}$, $\Z$ the integers, $[n]:=\{1,\dots,n\}$,
$\frac1a\Z:=\{z/a:\ z\in\Z\}$ for an integer $a\ge1$, and $\max\emptyset:=-\infty$.
\item $\F$ is a finite field, $\F^*:=\F\setminus\{0\}$, and $\E$ a finite extension of $\F$.
\item For $0\le i\le m$, a multi-index at level $i$ is $\alpha=(\alpha_{i+1},\dots,\alpha_m)\in\N^{m-i}$,
$x^\alpha:=x_{i+1}^{\alpha_{i+1}}\cdots x_m^{\alpha_m}$, and for a weight vector $w\in\N^m$,
$\langle w,\alpha\rangle:=\sum_{l=i+1}^mw_l\alpha_l$. $e_l\in\N^m$ is the $l$-th unit vector and
$\mathbf 1:=(1,\dots,1)$.
\item For a commutative $\F$-algebra $R$, every $q\in R[x_{i+1},\dots,x_m]$ is uniquely
$\sum_\alpha q_\alpha x^\alpha$ with $q_\alpha\in R$ (finitely many nonzero), and
$\deg_{x_l}q:=\max\{\alpha_l:\ q_\alpha\ne0\}$.
\item $\F[T]_{\le e}:=\{g:\deg g\le e\}$ and $\F[T]_{<\lambda}:=\{g:\deg g<\lambda\}$, with $\deg0:=-\infty$.
\item A word on a finite set $L$ is a map $u:L\to A$; $P|_L$ is $x\mapsto P(x)$.
\item A weight function on $L$ is $\mu:L\to[0,1]$; $\mu(A):=\frac1{|L|}\sum_{x\in A}\mu(x)$ for
$A\subseteq L$; $\mu$ has common denominator $M\ge1$ if $\mu(x)\in\frac1M\Z$ for all $x$.
\end{itemize}

\begin{definition}[Field and domains]\label{def:domains}
$\F$ is a finite field; $k\ge 2$, $m\ge 1$ and $N$ are integers with $k^m\mid N$ and
$N\mid |\F|-1$. $L_0\subset\F^*$ is a coset of the subgroup of order $N$, and
$L_i:=\{x^k: x\in L_{i-1}\}$ for $1\le i\le m$, so $|L_i|=N/k^i$. The \emph{fibre} of
$y\in L_i$ is $\fib(y):=\{x\in L_{i-1}: x^k=y\}$; it equals $x_0\Omega_k$ for any
$x_0\in\fib(y)$, where $\Omega_k:=\{\zeta\in\F:\ \zeta^k=1\}$ is the group of $k$-th roots of unity, and $|\fib(y)|=k$.
\end{definition}

Since $k\mid N\mid|\F|-1$, the cyclic group $\F^*$ has an element $\omega$ of order $k$, so
$\Omega_k$ is as in \cref{lem:rootsum}; and $k\ne0$ in $\F$, because
$\operatorname{char}\F\nmid|\F|-1$. Consequently $k^i|L_i|=N$ for $0\le i\le m$, the rate of
$\RS[\F,L_i,k^{m-i}]$ is $\rho:=k^m/N=k^{m-i}/|L_i|$ for every $i$, and $|L_i|\ge k^{m-i}$.

\begin{definition}[Polynomials]\label{def:poly}
$R$ denotes a commutative $\F$-algebra (in practice $R=\F$ or $R=\F[T]$). For
$1\le i\le m+1$, $R[x_i,\dots,x_m]_{<k}:=\{q\in R[x_i,\dots,x_m]:\ \deg_{x_l}q<k\text{ for }i\le l\le m\}$.
\end{definition}

\begin{definition}[Kronecker map]\label{def:kron}
For $0\le i\le m$ and $q\in R[x_{i+1},\dots,x_m]_{<k}$,
\[
U^{(i)}_q(Y):=q\big(Y,Y^{k},\dots,Y^{k^{m-i-1}}\big)\in R[Y],\qquad U_f:=U^{(0)}_f .
\]
For $i=m$, $q$ is a constant and $U^{(m)}_q=q$.
\end{definition}

\begin{definition}[Fold]\label{def:fold}
For $u:L_{i-1}\to R$ and $y\in L_i$, the \emph{components} of $u$ are
\[
u_j(y):=\frac1k\sum_{x\in\fib(y)}x^{-j}\,u(x)\qquad(0\le j<k),
\]
and for $c\in R$ the \emph{fold} is $\Fold_c(u):L_i\to R$,
$\Fold_c(u)(y):=\sum_{j<k}c^j\,u_j(y)$.
\end{definition}

\begin{definition}[Weighted and graded degree]\label{def:wdeg}
Let $w\in\N^m$. For $f=\sum_\alpha f_\alpha x^\alpha\in\F[x_1,\dots,x_m]$,
$\deg_w f:=\max\{\langle w,\alpha\rangle:\ f_\alpha\ne0\}$. For $0\le i\le m$ and
$q=\sum_{a,\alpha}q_{a,\alpha}T^ax^\alpha\in\F[T][x_{i+1},\dots,x_m]$ ($q_{a,\alpha}\in\F$),
$\gdeg_w q:=\max\{a+\langle w,\alpha\rangle:\ q_{a,\alpha}\ne0\}$.
\end{definition}

\begin{definition}[Codes]\label{def:codes}
$\RS[\F,L,d]:=\{P|_L:\deg P<d\}$. The $t$-interleaved code $\RS[\F,L,d]^t$ is the set
of words $L\to\F^t$ whose $t$ coordinates lie in $\RS[\F,L,d]$; we identify
$\F[T]_{<t}$ with $\F^t$ through coefficients. The \emph{graded codes} are
\[
\Gc_i^{w,D}:=\big\{\,U^{(i)}_q\big|_{L_i}\ :\ q\in\F[T][x_{i+1},\dots,x_m]_{<k},\ \gdeg_w q\le D\,\big\},
\]
and the \emph{degree code} $C_{w,D}$ is the set of $T$-free elements of $\Gc_0^{w,D}$.
\end{definition}

\begin{definition}[Weighted agreement]\label{def:agree}
For words $u,v$ on $L$ and $\mu:L\to[0,1]$,
$\agree_\mu(u,v):=\frac1{|L|}\sum_{x:\,u(x)=v(x)}\mu(x)$ (as in BCIKS20, \S7). For a set
$C$ of words, $\agree_\mu(u,C):=\max_{v\in C}\agree_\mu(u,v)$. Without a subscript,
$\mu\equiv1$: $\agree(u,v)=1-\Delta(u,v)$.
\end{definition}

Put $\Delta(u,v):=1-\agree(u,v)$; then $\Delta(u,w)\ge\Delta(v,w)-\Delta(v,u)$, because
$\{x:u(x)=w(x)\}\supseteq\{x:v(x)=w(x)\}\cap\{x:v(x)=u(x)\}$.

\paragraph{The caps.} For $w\in\N^m$ and $D\in\N$ put $d_0:=0$ and
$d_i:=\min\big(D,(k-1)\sum_{l\le i}w_l\big)$ for $1\le i\le m$; these are the formats of
\cref{def:protocol}.
